\chapter{Architecture et pipeline RAG}

% Partie technique : décrire le pipeline complet du RAG mis en place.

\section{Vue d'ensemble du pipeline}
% Schéma global :
% Documents Legifrance -> chunking -> embeddings -> Chroma
% Question -> embeddings -> retrieval -> reranking -> prompt -> LLM -> réponse
% Idéalement insérer une figure d'architecture.

\section{Chunking et indexation}
% - Stratégie de découpage (\texttt{CHUNK\_SIZE}, \texttt{CHUNK\_OVERLAP}).
% - Choix des paramètres et justification.
% - Embeddings : modèle utilisé via l'API UTC.
% - Stockage dans Chroma (cf. \texttt{rag\_pi/vectorstore.py}).

\section{Récupération et reranking}
% - Recherche par similarité dans Chroma (\texttt{RETRIEVER\_K}).
% - Reranking explicite : similarité vectorielle + bonus d'overlap lexical
%   (\texttt{RERANK\_OVERLAP\_WEIGHT}).
% - Seuil de pertinence (\texttt{MIN\_RELEVANCE\_SCORE}) pour filtrer les
%   documents peu pertinents.
% - Cf. \texttt{rag\_pi/rag.py}.

\section{Génération de la réponse}
% - Prompt RAG dédié au domaine juridique (cf. \texttt{rag\_pi/prompts.py}).
% - Contraintes imposées au modèle : citations d'articles, refus si pas
%   d'information suffisante.
% - Modèle utilisé (variable \texttt{LLM\_MODEL}) et endpoint UTC.

\section{Choix techniques principaux}
% Tableau ou liste des bibliothèques :
% LangChain, LangChain Google GenAI / OpenAI-compatible, Chroma, Streamlit.
